\documentclass[12pt, a4paper]{article}
\usepackage{xeCJK} % Support for Chinese characters
\setCJKmainfont{Microsoft JhengHei}[
    AutoFakeBold=2.5,
    AutoFakeItalic=true
]
\usepackage{geometry}
\geometry{left=2.5cm, right=2.5cm, top=2.5cm, bottom=2.5cm}
\usepackage{booktabs}
\usepackage{float}

\begin{document}

\noindent
\textbf{個人報告} \hfill 系級：資工二 \quad 學號：113590051 \quad 姓名：楊承諭

\vspace{1.5em}
\begin{center}
    {\Large \textbf{個人心得報告}}
\end{center}
\vspace{1em}

本次 Lab 08 主要學習以 VHDL 實作一簡易中央處理器（Simple CPU）的硬體設計，並在 FPGA 開發板上完成實體電路驗證。本次實驗設計的 CPU 架構主要由暫存器檔案（Register File）、內部資料匯流排（Data Bus）、指令譯碼器（Instruction Decoder）以及輸出解碼模組所組成，支援 \texttt{LOAD} 與 \texttt{MOVE} 兩大核心指令，並藉由有限狀態機（FSM）或同步循序邏輯調控暫存器資料的存取與傳輸。

在設計過程中，我深入理解了簡易 CPU 內部資料路徑（Data Path）的規劃與運作機制。本處理器內建四個八位元通用暫存器（$R0\sim R3$），暫存器的讀取採組合電路設計，透過 Rs 與 Rt 的選取線即時輸出暫存器內容；而寫入則採同步循序邏輯，在時脈正緣觸發時，依據指令碼（Opcode）將資料寫入對應的目標暫存器 Rs。指令格式設計上，我們將 16 位元指撥開關（SW[15:0]）進行欄位規劃，其中低 8 位元（SW[7:0]）作為立即值資料，SW[11:8] 作為 Opcode，SW[13:12] 與 SW[15:14] 則分別為目標與來源暫存器選取線。

本次實作中最關鍵的挑戰在於時序控制與按鍵防彈跳的處理。我們採用 KEY[0] 作為手動時脈輸入，由於按鍵預設為低電平有效（active-low），為了讓 CPU 能在按鍵按下的瞬間精確執行指令，我設計了 \texttt{clk} $\leftarrow$ \texttt{not KEY(0)} 的反相電路，將其轉換為正緣觸發。在硬體測試階段，當指撥開關切換至 \texttt{LOAD R1, \#0x55} 時，按壓 KEY[0] 後，內部匯流排與 R1 的數值立即同步變更，並由七段顯示器 HEX0--HEX5 正確呈現匯流排與 Rs/Rt 的值。這次的專案製作與實機測試，不僅強化了我對處理器架構與控制單元的設計能力，更讓我對數位系統的整合開發與實體偵錯有了更深層的掌握。

\vspace{2em}
\noindent
\textbf{工作分配表}
\vspace{0.5em}

\begin{table}[H]
    \centering
    \begin{tabular}{llcl}
        \toprule
        \textbf{學號} & \textbf{姓名} & \textbf{比重} & \textbf{工作項目} \\ \midrule
        113590051 & 楊承諭 & 65\% & 小組專案製作與測試 \\
        113590017 & 黃楷程 & 35\% & 小組報告製作 \\ \bottomrule
    \end{tabular}
\end{table}

\end{document}
